\chapter{Glossary}
\label{ch:glossary}

Every technical term used in the paper and in this compendium, in alphabetical order,
with one to three sentences and the chapter in which it is defined. Symbols follow the
paper's notation (\Cref{ch:conventions}).

\begingroup\sloppy
\begin{description}[leftmargin=1.6em,itemsep=0.25em,font=\normalfont\bfseries]

\item[Adam] A first-order stochastic optimiser that scales each parameter's update by
running estimates of the gradient's mean and magnitude; the optimiser of every network
arm. \Cref{ch:neural-baselines}.

\item[Affine layer] A map $\vect{h}\mapsto\mat{W}\vect{h}+\vect{b}$; the basic block of every
network, with $mn+m$ parameters for an $m\times n$ weight. \Cref{ch:neural-baselines}.

\item[Amendment, pre-data] A dated entry in a registry file recording a change to a
registered design or analysis made before any confirmatory cell has produced a number;
the instrument by which a registration is corrected without losing its guarantee.
\Cref{ch:protocol}.

\item[Amplitude (of a feature)] The size of a Hermite moment of degree $m$, attenuated by
the channel as $e^{-mt}$; the first of the three clocks. \Cref{ch:attenuation-theorems}.

\item[AR(1) chain] A first-order autoregression $a_i=\corr a_{i-1}+\varepsilon_i$ with
i.i.d.\ innovations; the paper's data model, with $a_1\sim\normal01$ and
$\ivar=1-\corr^2$. \Cref{ch:markov-chains}.

\item[AR(2)] A second-order autoregression $a_i=\corr_1a_{i-1}+\corr_2a_{i-2}+\varepsilon_i$;
one of the four misspecification mechanisms, parameterised by $\corr_2$ at fixed lag-one
correlation. \Cref{ch:misspecification-edgeworth}.

\item[Assumption (exponential moment)] $\E e^{\lambda|\va|}<\infty$ for some $\lambda>0$,
the hypothesis of the paper's Theorem~3; every family used except Student-$t$ satisfies
it. \Cref{ch:attenuation-theorems}.

\item[Asymptotic relative efficiency] The limiting ratio of the Cram\'er--Rao variance to
an estimator's variance; one for an efficient estimator, below one otherwise.
\Cref{ch:quasi-likelihood-moments}.

\item[At cap] A cell whose validation-selected checkpoint is the last rung of its ladder,
so the reported error is an upper bound on what a longer ladder would give.
\Cref{ch:model-selection}.

\item[Attention (self-attention)] The transformer's operation in which every site forms a
softmax-weighted average of every other site's representation with weights it computes
from the data. \Cref{ch:neural-baselines}.

\item[Attenuation, Hermite] The exact identity $\E_{\Pt}[H_\nu(\vx_t)]=e^{-|\nu|t}\E_{\Pz}[H_\nu(\va)]$,
the paper's Theorem~2. \Cref{ch:attenuation-theorems}.

\item[Audit script] \code{tools/audit\_claims.py}, an independent recomputation of every
headline number from the frozen CSVs with plain code, compared with the macros.
\Cref{ch:protocol}.

\item[Backward message] $\rmsg_i$, the function of $a_i$ carrying the evidence from sites
$i+1,\dots,\nsites$; the right-to-left sweep of sum-product. \Cref{ch:factor-graphs-bp}.

\item[Belief (site marginal)] The posterior density of one site,
$\propto\lmsg_i\like\rmsg_i$; its mean is $\pmeansite{i}$. \Cref{ch:factor-graphs-bp}.

\item[Belief propagation] Sum-product message passing on a factor graph; exact on a tree,
hence exact for the paper's posterior. \Cref{ch:factor-graphs-bp}.

\item[Bias (of an estimator)] $\E\hat\theta-\theta$; in misspecification, the error that
remains as $\nseq\to\infty$. \Cref{ch:fisher-cramer-rao,ch:misspecification-edgeworth}.

\item[Bi-GRU] A bidirectional gated recurrent network; one of the six baselines.
\Cref{ch:neural-baselines}.

\item[Bimodal innovation] The symmetric two-component Gaussian mixture innovation (separation
parameter $0.6$), excess kurtosis $-0.72$. \Cref{ch:markov-chains}.

\item[Bootstrap, percentile] Resampling seeds with replacement and reporting the $2.5$th
and $97.5$th percentiles of the resampled statistic; the paper's interval for paired
ratios. \Cref{ch:experimental-statistics}.

\item[Break-even] The misspecification strength at which the paired log-ratio of the
network's error to EM--BP's changes sign; reported as the bracketing grid points.
\Cref{ch:misspecification-edgeworth}.

\item[Bulk region] Sites $5$ to $\nsites-4$; the region used for checkpoint selection.
\Cref{ch:model-selection}.

\item[Bundle] A fixed set of chains with their noised views at every level and the
reference posterior means under the true kernel; training, validation and test bundles
are drawn from disjoint keyed streams. \Cref{ch:model-selection}.

\item[Capacity (mixture)] The number of components $\ncomp$ of the mixture-innovation
kernel. \Cref{ch:em-through-noise}.

\item[Cell] One unit of work of a driver (one seed at one size, say), written to its own
file so that a run is resumable. \Cref{ch:protocol}.

\item[Checkpoint ladder] The fixed list of budgets (EM iterations or gradient steps) at
which a fit is saved; the selected checkpoint is chosen afterwards among its rungs.
\Cref{ch:model-selection}.

\item[Clean chain] The unnoised sequence $\va\sim\Pz$. \Cref{ch:conventions}.

\item[Closure gap] $\E_{\Pt}\lVert\score-\Sgauss\rVert^2$, the squared distance between
the true score and the covariance twin's; $\tfrac16\Frob{\kfour}^2e^{-8t}$ to leading order.
\Cref{ch:attenuation-theorems}.

\item[Coloured probes] The $2r+1$ deterministic probe vectors that give the exact trace of
a Jacobian of bandwidth $r$; the divergence estimator of the window head.
\Cref{ch:neural-baselines}.

\item[Common random numbers] Keying every random stream by a tag so that compared arms see
the same draws and their difference is not Monte Carlo noise. \Cref{ch:protocol}.

\item[Conditional expectation] $\E[Y\mid X]$, the best mean-square predictor of $Y$ from
$X$; the posterior mean is the case $Y=\va$, $X=\vx_t$. \Cref{ch:conventions}.

\item[Confirmatory seeds] Seeds $0$--$15$ (eight in the sweeps), used only for reported
numbers, never for selection. \Cref{ch:model-selection}.

\item[Consistency] Convergence of an estimator to the true (or pseudo-true) parameter as
$\nseq\to\infty$. \Cref{ch:fisher-cramer-rao,ch:quasi-likelihood-moments}.

\item[Convolution, one-dimensional] A weight-shared linear map of a signal whose output
at $i$ depends on inputs within a fixed radius; zero-padded at the ends.
\Cref{ch:neural-baselines}.

\item[Correlation coordinate] A parameter that moves the covariance $\Sig_0$, such as
$\corr$; a degree-two coordinate. \Cref{ch:attenuation-theorems}.

\item[Covariance twin] The Gaussian chain with the same $\corr$ and $\ivar$; its score is
$\Sgauss=-\Sig_t^{-1}\vx$. \Cref{ch:gaussian-closure-lmmse}.

\item[Cram\'er--Rao bound] $\Var\hat\theta\ge1/(\nseq I(\theta))$ for an unbiased estimator;
the reference curve of the learning figure. \Cref{ch:fisher-cramer-rao}.

\item[Crossover time] $\tstar=\tfrac12\log\lmax(\Sig_0)$, where $e^{2t}$ passes the
largest eigenvalue of the clean covariance; the scale beyond which the asymptotic rates
apply. \Cref{ch:attenuation-theorems}.

\item[Cumulant] The coefficients of the log of the moment generating function; additive
over independent summands and homogeneous of degree $n$. \Cref{ch:moments-cumulants}.

\item[Cumulant generating function] $\log\E e^{s\cdot\va}$, whose Taylor coefficients are
the joint cumulants. \Cref{ch:moments-cumulants}.

\item[Damped cumulant generating function] $\Lambda(s,b)=\log\E_{\Pz}\exp(s\cdot\va-b|\va|^2/2)$,
the function whose grading gives the paper's Theorem~3. \Cref{ch:attenuation-theorems}.

\item[Delta method] The first-order propagation of an estimator's variance through a
smooth function. \Cref{ch:experimental-statistics}.

\item[Denoiser] A map from a noised sequence and a level to an estimate of the posterior
mean; every network arm is one. \Cref{ch:neural-baselines}.

\item[Denoising score matching] Regressing the network on the clean chain (or the noise)
with fresh corruptions; its minimiser is the posterior mean, hence the score.
\Cref{ch:score-matching,ch:neural-baselines}.

\item[Development seeds] Seeds $10000$--$10003$, used only for architecture and
hyperparameter choices. \Cref{ch:model-selection}.

\item[Dilation] The spacing $d$ between the taps of a convolution; a dilated stack widens
its receptive radius by $d$ per layer. \Cref{ch:neural-baselines}.

\item[Divergence] $\nabla\cdot s=\sum_i\partial s_i/\partial x_i$, the trace of the score
map's Jacobian; the term implicit score matching must estimate. \Cref{ch:neural-baselines}.

\item[Drift (grid)] The relative change of the posterior-mean field between two grids on
the same chains; the quadrature convergence diagnostic. \Cref{ch:grid-quadrature}.

\item[Edgeworth correction] The expansion of a law around its covariance-matched Gaussian,
whose first term is $\tfrac1{24}\kfour$ contracted with the multivariate Hermite
polynomials. \Cref{ch:misspecification-edgeworth}.

\item[Efficient information] The Fisher information of a coordinate after the nuisance
coordinates are projected out. \Cref{ch:fisher-cramer-rao}.

\item[EM (expectation--maximisation)] Alternating an E-step (posterior expectations under
the current parameters) and an M-step (maximisation of the expected complete-data
log-likelihood). \Cref{ch:maximum-likelihood-em}.

\item[EM--BP] The paper's structured estimator: EM whose E-step is the tree-exact
recursion on the grid and whose kernel is the mixture-innovation AR(1) kernel.
\Cref{ch:em-through-noise}.

\item[Eps parameterisation] The network predicts the standardised noise $\vz$; the mean
is $(\vx-\sqrt{\Dt}\hat\varepsilon)/e^{-t}$. \Cref{ch:neural-baselines}.

\item[Estimand] The quantity a number estimates; the paper's is the schedule-pooled
relative score error. \Cref{ch:experimental-statistics}.

\item[Evidence (marginal likelihood)] $\Pt(\vx)$, the probability of the noised view under
the model, returned by the forward sweep with its constants accumulated.
\Cref{ch:grid-quadrature}.

\item[Evidence rule] Selecting EM's rung by the held-out one-view log-evidence per edge;
the oracle-free rule of E20 and E14. \Cref{ch:model-selection}.

\item[Exceedance count] The expected number of sites of a bundle beyond the grid domain,
$\nseq\nsites P(|a|>\halfwidth)$; the right diagnostic for polynomial tails.
\Cref{ch:grid-quadrature}.

\item[Excess kurtosis] $\kfour/\ivar^2$, the fourth cumulant divided by the squared variance;
$+3$ for Laplace, $0$ for Gaussian. \Cref{ch:moments-cumulants}.

\item[Factor graph] A bipartite graph of variables and factors representing a product
density; the posterior's is the prior's chain. \Cref{ch:factor-graphs-bp}.

\item[Fair screen] A screen in which every candidate is trained under the same recipe,
budget, bundles, ladder and selection rule on development seeds disjoint from the
confirmatory ones. \Cref{ch:model-selection}.

\item[Fisher divergence] $\tfrac12\E_p\lVert s_\theta-s_p\rVert^2$, what score matching
minimises. \Cref{ch:score-matching,ch:model-selection}.

\item[Fisher information] $I(\theta)=\E[(\partial_\theta\log P)^2]$; per noised sequence,
$I_t(\theta)$. \Cref{ch:fisher-cramer-rao}.

\item[Fisher's identity] The gradient of the marginal log-likelihood equals the posterior
expectation of the complete-data score; on the grid it is read off the pairwise
marginals. \Cref{ch:maximum-likelihood-em}.

\item[Forward message] $\lmsg_i$, the function of $a_i$ carrying the evidence from sites
$1,\dots,i-1$. \Cref{ch:factor-graphs-bp}.

\item[Forward process] The Ornstein--Uhlenbeck noising $\vx_t=e^{-t}\va+\sqrt{\Dt}\vz$.
\Cref{ch:ou-channel-tweedie}.

\item[Fourier features] The vector $(\log t,\sin2^k\log t,\cos2^k\log t)_{k\le3}$ by which the
level enters every network. \Cref{ch:neural-baselines}.

\item[Frobenius norm] $\Frob{\mat{A}}^2=\sum A_{ij}^2$, extended to tensors over ordered
index tuples. \Cref{ch:conventions}.

\item[Frozen configuration] \code{FROZEN}, the world every experiment runs in: chain,
channel, grid, schedule, seeds. \Cref{ch:protocol}.

\item[Frozen output] An output directory whose files are hashed into its registry entry;
immutable thereafter. \Cref{ch:protocol}.

\item[Gaussian closure] Moment-projected sum-product, which replaces every message by the
Gaussian with its mean and variance; equals the LMMSE of the covariance twin.
\Cref{ch:gaussian-closure-lmmse}.

\item[Gaussian quasi-likelihood] The Gaussian log-likelihood written for data that are
not Gaussian; depends on the data only through the sample covariance.
\Cref{ch:quasi-likelihood-moments}.

\item[GELU] The activation $x\Phi(x)$, a smooth rectifier. \Cref{ch:neural-baselines}.

\item[Generalised EM] An EM whose M-step only increases the expected complete-data
log-likelihood (conditional maximisation); its ascent is monotone.
\Cref{ch:maximum-likelihood-em}.

\item[Generalised Gaussian family] Innovations $\propto\exp(-|\varepsilon/s_\beta|^\beta)$
at fixed variance, from Laplace ($\beta=1$) through Gaussian ($\beta=2$) to platykurtic
members. \Cref{ch:markov-chains}.

\item[Generator] A script under \code{tools/} that reads frozen outputs and writes the
macros, tables and figures of the paper. \Cref{ch:protocol}.

\item[Gradient clipping] Rescaling the gradient to a maximum norm before an optimiser
step; at $1$ in the recipe. \Cref{ch:neural-baselines}.

\item[Grid (quadrature)] The $\ngrid$ uniform points on $[-\halfwidth,\halfwidth]$ on which
messages are carried with trapezoidal weights. \Cref{ch:grid-quadrature}.

\item[GRU] Gated recurrent unit; a recurrence with reset and update gates.
\Cref{ch:neural-baselines}.

\item[H\'ajek--Le Cam] The local asymptotic minimax theorem: no estimator beats
$1/I^{\mathrm{eff}}$ locally, and maximum likelihood attains it.
\Cref{ch:fisher-cramer-rao}.

\item[Heavy tails] Innovation laws with polynomial tails (Student-$t$), outside the
exponential-moment assumption, with an $e^{-\nu t}$ term in the expected gap.
\Cref{ch:attenuation-theorems}.

\item[Held-out] Data not used to fit; the validation bundle for selection, the test
bundle for scoring. \Cref{ch:model-selection}.

\item[Hermite moment] $\E[H_\nu(\vx)]$ with $H_\nu$ a product of Hermite polynomials over
sites; the quantity the channel attenuates diagonally. \Cref{ch:gaussian-hermite}.

\item[Hermite polynomial] The probabilists' $\Hen{n}$, orthogonal under the standard
Gaussian with $\E\Hen{n}^2=n!$. \Cref{ch:gaussian-hermite}.

\item[Hutchinson's estimator] $\vect{v}^\top\mat{J}\vect{v}$ with a standard normal probe,
an unbiased but noisy estimate of $\tr\mat{J}$. \Cref{ch:neural-baselines}.

\item[Hyv\"arinen identity] $\E_p[\tfrac12\lVert s\rVert^2+\nabla\cdot s]=\tfrac12\E_p\lVert s-s_p\rVert^2+\text{const}$,
by integration by parts. \Cref{ch:score-matching,ch:model-selection}.

\item[Identifiability] Distinct parameters give distinct laws; $(\corr,\ivar)$ is
identifiable from $\Sig_0$, hence from the noised covariance.
\Cref{ch:quasi-likelihood-moments}.

\item[Implicit score matching] Minimising Hyv\"arinen's objective on noisy samples only;
the loss of interface~C. \Cref{ch:score-matching,ch:neural-baselines}.

\item[Information interface] What a learner may see: A (clean chains, fresh corruptions),
B (clean chains, one fixed corruption), C (one noised view per sequence and nothing
else). \Cref{ch:model-selection}.

\item[Innovation] The i.i.d.\ increment $\varepsilon_i$ of the chain, with density $\innov$
of mean zero and variance $\ivar$. \Cref{ch:markov-chains}.

\item[Isotropic score] $-\vx$, the score of $\normal0{\Id}$; the reference of a model that
learned nothing beyond the unit variance. \Cref{ch:gaussian-closure-lmmse}.

\item[Isserlis' theorem] Gaussian moments are sums over pairings of covariances (also
Wick's theorem). \Cref{ch:gaussian-hermite}.

\item[k-statistic] Fisher's exactly unbiased sample cumulant. \Cref{ch:quasi-likelihood-moments}.

\item[Kalman smoother] The linear-Gaussian smoother; the Gaussian part of the paper's
score. \Cref{ch:gaussian-closure-lmmse}.

\item[Kernel (transition)] $\kernel(a'\mid a)=\innov(a'-\corr a)$; $\kernelth$ when
parameterised. \Cref{ch:markov-chains}.

\item[Kullback--Leibler divergence] $\KL(p\|q)=\E_p\log(p/q)\ge0$, zero only at $p=q$.
\Cref{ch:conventions}.

\item[Kullback--Leibler projection] The member of a family closest to a law in
$\KL(\text{law}\|\cdot)$; the limit of a misspecified maximum-likelihood estimator.
\Cref{ch:misspecification-edgeworth}.

\item[Lag-one moment estimator] $\hat\corr=e^{2t}\,\overline{x_ix_{i+1}}$, unbiased for
$\corr$ by Hermite attenuation. \Cref{ch:quasi-likelihood-moments}.

\item[Laplace innovation] The double-exponential density, excess kurtosis $3$; the
frozen family. \Cref{ch:markov-chains}.

\item[Layer normalisation] Rescaling a representation to zero mean and unit variance
across channels with a learned gain and shift. \Cref{ch:neural-baselines}.

\item[Length-agnostic] An architecture that refers to no fixed $\nsites$ and can be
applied at any length. \Cref{ch:neural-baselines}.

\item[LMMSE] The linear minimum-mean-square estimator $e^{-t}\Sig_0\Sig_t^{-1}\vx$ of
$\va$ from $\vx_t$; the best linear map. \Cref{ch:gaussian-closure-lmmse}.

\item[LMMSE floor] The relative score error of the LMMSE with the true covariance; the
bound of every second-order estimator. \Cref{ch:gaussian-closure-lmmse}.

\item[Locality floor] The error of the best estimator that sees only a window of radius
$r$, the exact posterior mean given the window. \Cref{ch:attenuation-theorems}.

\item[Log-evidence per edge] The held-out marginal log-likelihood divided by the number
of transitions; the evidence rule's statistic. \Cref{ch:model-selection}.

\item[Log-spaced schedule] The twelve levels $0.05\cdot60^{(k-1)/11}$. \Cref{ch:protocol}.

\item[M-estimator] The solution of an estimating equation $\sum_\mu\psi(\vx^\mu;\theta)=0$;
its variance is the sandwich. \Cref{ch:quasi-likelihood-moments}.

\item[M-step] The maximisation step of EM; for a Gaussian kernel the belief-weighted
Yule--Walker pair. \Cref{ch:maximum-likelihood-em}.

\item[Marginal likelihood] See evidence. \Cref{ch:maximum-likelihood-em}.

\item[Markov chain] A sequence in which each site depends on the past only through its
predecessor. \Cref{ch:markov-chains}.

\item[Matched interface] Interface~C, EM--BP's own; the comparison at matched information.
\Cref{ch:model-selection}.

\item[Mehler's identity] $\E_z\Hen{n}(\alpha a+\sqrt{1-\alpha^2}z)=\alpha^n\Hen{n}(a)$; the
proof of Hermite attenuation. \Cref{ch:gaussian-hermite}.

\item[Method of moments] Solving sample moments equal to model moments for the parameter.
\Cref{ch:quasi-likelihood-moments}.

\item[Misspecification] Fitting a family that does not contain the truth.
\Cref{ch:misspecification-edgeworth}.

\item[Mixture-innovation kernel] $\kernelth(a'\mid a)=\sum_c\pi_c\normal{a'-\corr a-\nu_c}{s_c^2}$,
the paper's parametric kernel. \Cref{ch:em-through-noise}.

\item[MLP] The fully connected network; the baseline with no weight sharing.
\Cref{ch:neural-baselines}.

\item[Moment] $\E[X^n]$ or its multivariate version. \Cref{ch:moments-cumulants}.

\item[Multi-index] $\nu\in\mathbb N^{\nsites}$ with $|\nu|=\sum\nu_i$, labelling a product
of Hermite polynomials. \Cref{ch:gaussian-hermite}.

\item[Multivariate Hermite polynomial] $H^\Sig_{ijkl}=\partial_i\partial_j\partial_k\partial_l\phi_\Sig/\phi_\Sig$,
the tensor version of $\Hen{4}$ for a Gaussian with covariance $\Sig$.
\Cref{ch:misspecification-edgeworth}.

\item[Noised view] One draw $\vx=e^{-t}\va+\sqrt{\Dt}\vz$ at one level.
\Cref{ch:ou-channel-tweedie}.

\item[One-view protocol] Each training sequence contributes one noised view at one level,
the sequences split evenly across the twelve levels. \Cref{ch:em-through-noise}.

\item[Oracle] A quantity that uses the truth (the true kernel's posterior means, the true
score); an oracle rule is a selection rule that uses one. \Cref{ch:model-selection}.

\item[Oracle selection rule] Choosing the rung by the validation risk against the true
kernel's posterior means on the bulk. \Cref{ch:model-selection}.

\item[Ornstein--Uhlenbeck channel] The coordinatewise noising process with decay $e^{-t}$
and variance $\Dt=1-e^{-2t}$, which keeps the unit variance. \Cref{ch:ou-channel-tweedie}.

\item[Outer-product information estimate] Estimating Fisher information by the sample
covariance of per-sequence scores. \Cref{ch:fisher-cramer-rao}.

\item[Overfull box] A line LaTeX could not fit in its width; a gate of the manuscript
check for tables and figures. \Cref{ch:protocol}.

\item[Paired comparison] A difference or ratio formed within seed on identical bundles
before averaging over seeds. \Cref{ch:experimental-statistics}.

\item[Parameter count] The number of trainable numbers of a network, fixed by its
architecture. \Cref{ch:neural-baselines}.

\item[Pointwise MLP] Affine maps and activations applied at each site separately
(width-one convolutions). \Cref{ch:neural-baselines}.

\item[Pooled relative score error] $\bigl(\sum_t\sum_{j,i}(\hat s-s)^2/\sum_t\sum_{j,i}s^2\bigr)^{1/2}$,
the paper's estimand. \Cref{ch:experimental-statistics}.

\item[Posterior mean] $\pmean{\va}=\E[\va\mid\vx_t=\vx]$; the score is
$-(\vx-e^{-t}\pmean{\va})/\Dt$. \Cref{ch:ou-channel-tweedie}.

\item[Pre-registration] Writing hypothesis, design, seeds and analysis before the data
exist. \Cref{ch:protocol}.

\item[Provenance record] \code{params\_<part>.json}: commit, dirty list, archive digest,
host, resolved settings. \Cref{ch:protocol}.

\item[Pseudo-true parameter] The minimiser of $\KL(P^\star\|Q_\theta)$; where a
misspecified estimator converges. \Cref{ch:misspecification-edgeworth}.

\item[Quadrature, trapezoidal] Approximating an integral by a weighted sum on a uniform
grid; second order in the spacing for smooth integrands. \Cref{ch:grid-quadrature}.

\item[Quasi-MLE] The maximiser of a quasi-likelihood; consistent for $(\corr,\ivar)$ here.
\Cref{ch:quasi-likelihood-moments}.

\item[Quasi-score] The derivative of the quasi-likelihood; a moment equation in the
sample covariance. \Cref{ch:quasi-likelihood-moments}.

\item[Receptive radius] The $r$ such that a network's output at site $i$ depends on inputs
within $|j-i|\le r$ only. \Cref{ch:neural-baselines}.

\item[Registry] \code{experiments/registry/<id>.json}: hypothesis, design, analysis, seeds,
status, amendments, frozen digests. \Cref{ch:protocol}.

\item[Relative-position bias] A learned attention bias depending on the clipped offset
$i-j$; what makes the transformer length-agnostic. \Cref{ch:neural-baselines}.

\item[Residual connection] $\vect{h}\mapsto\vect{h}+g(\vect{h})$. \Cref{ch:neural-baselines}.

\item[Reverse SDE] The time-reversed diffusion driven by the score, used to generate
samples. \Cref{ch:reverse-sde-generation}.

\item[RMSE] Root mean squared error over seeds of a fitted coordinate.
\Cref{ch:experimental-statistics}.

\item[Rung] One entry of a checkpoint ladder. \Cref{ch:model-selection}.

\item[Sample covariance] $\mat{S}=\nseq^{-1}\sum_\mu\vx^\mu\vx^{\mu\top}$.
\Cref{ch:quasi-likelihood-moments}.

\item[Sandwich covariance] $\mat{J}^{-1}\mat{V}\mat{J}^{-1}$, the asymptotic variance of an
M-estimator under a possibly wrong model. \Cref{ch:quasi-likelihood-moments}.

\item[Schedule] The twelve levels over which errors are pooled. \Cref{ch:protocol}.

\item[Score] $\score(\vx,t)=\nabla_{\vx}\log\Pt(\vx)$. \Cref{ch:ou-channel-tweedie}.

\item[Score matching] Fitting a score by minimising the Fisher divergence, implicitly or
by denoising. \Cref{ch:score-matching}.

\item[Screen] The selection of a configuration per architecture family on development
seeds. \Cref{ch:model-selection}.

\item[Seed] The integer that determines a random stream; the inferential unit of the
paper's comparisons. \Cref{ch:experimental-statistics}.

\item[SHA-256] A cryptographic hash producing a $256$-bit digest; a changed byte changes
the digest. \Cref{ch:protocol}.

\item[Shape coordinate] A parameter that leaves $\Sig_0$ fixed and moves $\kfour$; a
degree-four coordinate, such as the generalised-Gaussian exponent at fixed variance.
\Cref{ch:attenuation-theorems}.

\item[Site] One coordinate $i$ of the chain. \Cref{ch:conventions}.

\item[Sliced Wasserstein distance] A distance between samples averaged over random
one-dimensional projections; a column of the generation table.
\Cref{ch:reverse-sde-generation}.

\item[SLURM] The cluster scheduler that runs the registered experiments. \Cref{ch:protocol}.

\item[Sub-chain] A window of consecutive sites treated as a chain with the exact marginal
of its first site as initial factor; the locality floor's computation.
\Cref{ch:attenuation-theorems}.

\item[Sum-product] The message-passing algorithm computing marginals on a factor graph.
\Cref{ch:factor-graphs-bp}.

\item[TCN] The dilated temporal convolutional network, receptive radius $15$.
\Cref{ch:neural-baselines}.

\item[Test bundle] The $2048$ held-out chains on which a selected fit is scored.
\Cref{ch:model-selection}.

\item[Three clocks] Amplitude $e^{-mt}$, information $e^{-2mt}$, samples $e^{2mt}$ for a
feature first visible at Hermite degree $m$. \Cref{ch:attenuation-theorems}.

\item[Time features] See Fourier features. \Cref{ch:neural-baselines}.

\item[Transformer] The attention-based encoder with relative-position bias.
\Cref{ch:neural-baselines}.

\item[Truncation (domain)] The error from carrying messages on a finite domain; bounded
by the exceedance count, not by the grid spacing. \Cref{ch:grid-quadrature}.

\item[Tweedie's identity] $\score(\vx,t)=-(\vx-e^{-t}\pmean{\va})/\Dt$.
\Cref{ch:ou-channel-tweedie}.

\item[Typed literal] A decimal typed into the manuscript rather than emitted by a
generator; forbidden unless listed with a justification. \Cref{ch:protocol}.

\item[U-Net] The two-resolution convolutional network with a skip connection.
\Cref{ch:neural-baselines}.

\item[Unary factor] $\like(x_i\mid a_i)\propto\exp[-(x_i-e^{-t}a_i)^2/2\Dt]$, the
likelihood factor attached to each site. \Cref{ch:factor-graphs-bp}.

\item[Uniform innovation] The uniform density at unit variance, excess kurtosis $-1.2$.
\Cref{ch:markov-chains}.

\item[Validation bundle] The $1024$ held-out chains on which the checkpoint is chosen.
\Cref{ch:model-selection}.

\item[Validation risk] The pooled relative score error on the validation bundle's bulk
against the true kernel's posterior means. \Cref{ch:model-selection}.

\item[Variance-matched family] Innovation laws with the same mean and variance, hence the
same $\Sig_0$, differing only beyond second order. \Cref{ch:markov-chains}.

\item[Weight sharing] Using the same parameters at every site, as a convolution or a
recurrence does. \Cref{ch:neural-baselines}.

\item[Wick product] ${:}x_ix_jx_k{:}$, the Hermite (orthogonalised) monomial; Wick's
theorem gives its Gaussian second moments by pairings. \Cref{ch:gaussian-hermite}.

\item[Window head] One convolution of width $2r+1$ followed by a pointwise MLP; the
smallest local baseline. \Cref{ch:neural-baselines}.

\item[Winner's curse] The downward bias of the minimum of several noisy measured risks;
the reason selection and reporting use disjoint seeds. \Cref{ch:model-selection}.

\item[x0 parameterisation] The network predicts the posterior mean directly.
\Cref{ch:neural-baselines}.

\item[Yule--Walker equations] The linear relations between an autoregression's
coefficients and its autocovariances; the Gaussian M-step and the AR(2) covariance.
\Cref{ch:markov-chains,ch:maximum-likelihood-em}.

\item[Zero padding] Treating inputs outside $\{1,\dots,\nsites\}$ as zero in a convolution;
the source of the window head's boundary error. \Cref{ch:neural-baselines}.

\end{description}
\endgroup
